% !TeX root = ../../Book4.tex
% 纲要标题：### 15.2 谱序列的构造
% 纲要标签：b4:sec:1402
% 纲要位置：计划纲要.md，第 2833 行。

\section{谱序列的构造}
\label{b4:sec:1402}

第零页只看伴随分次，后续微分则测量代表元还能向深层提升多少次。本节先从代表元直接构造各页，再证明每一新页由前一页取同调得到。

% 纲要内容：
%
% * E_0、E_1 及微分
% * 页间同调关系
%

\subsection{\texorpdfstring{\(E_0\)、\(E_1\) 与微分}{E0、E1 与微分}}
\label{b4:subsec:1402:01}

设 \(\mathcal A\) 为交换范畴，\((C,F)\) 为 \(\mathcal A\) 中的递减滤过上链复形。第零页为
\[
 E_0^{p,q}=F^pC^{p+q}/F^{p+1}C^{p+q},\qquad \mathrm{d}_0[x]=[\mathrm{d}x].
\]
微分保持滤过，保证它良定义。因此
\[
 E_1^{p,q}=H^{p+q}(F^pC/F^{p+1}C).
\]
若 \(x\) 代表此页的类，便有 \(\mathrm{d}x\in F^{p+1}\)。再取 \(\mathrm{d}x\) 在下一层的类，就得到 \(\mathrm{d}_1[x]\)。

\ProseParagraphBreak
例如，取零次和一次的两项复形 \(C=(\mathbb Z\xrightarrow{2}\mathbb Z)\)，并令 \(F^pC=C\)（\(p\le0\)）、\(F^1C=(0\to\mathbb Z)\)、\(F^pC=0\)（\(p\ge2\)）。乘 \(2\) 把源的滤过层零送入靶的滤过层一，所以伴随分次上的 \(\mathrm{d}_0\) 为零。第零页与第一页都只有两个非零项，但第一页微分恢复了原复形中跨过一层滤过的映射：
\[
 E_1^{0,0}=\mathbb Z\xrightarrow{\ \mathrm{d}_1=\times2\ }
 E_1^{1,0}=\mathbb Z.
\]
于是第二页只剩 \(E_2^{1,0}=\mathbb Z/2\)，以后再无可能的非零微分。它与 \(H^0(C)=0\)、\(H^1(C)=\mathbb Z/2\) 相符：先逐层看时得到的两个自由群，通过后续微分才合成为原复形的挠同调。

\ProseParagraphBreak
为同时写出后面各页，令 \(n=p+q\)，并对 \(r\ge0\) 定义
\[
 Z_r^{p,q}=F^pC^n\cap \mathrm{d}^{-1}(F^{p+r}C^{n+1}).
\]
这里子对象的交通过两嵌入的拉回定义，\(\mathrm d^{-1}(F^{p+r}C^{n+1})\) 是沿微分的拉回原像；两个子对象的和是它们的直和映入 \(C^n\) 所得的像。每页的公式只涉及有限交、有限和及商对象。特别地 \(Z_0^{p,q}=F^pC^n\)。当 \(r\ge1\) 时，置
\begin{equation}\label{b4:eq:spectral-page-formula}
 E_r^{p,q}=
 \frac{Z_r^{p,q}}
 {Z_{r-1}^{p+1,q-1}+\mathrm{d}Z_{r-1}^{p-r+1,q+r-2}}.
\end{equation}
分母第一项容许更改较深一层的代表元，但其微分也要足够深；第二项容许减去已经可见的余边。两项确实包含于分子：第一项的微分位于 \(F^{p+r}\)，第二项位于 \(F^p\) 且再次微分为零。取 \(r=1\)，公式正是上面第零页的上同调。

\subsection{页间同调关系}
\label{b4:subsec:1402:02}

\begin{theorem}\label{b4:thm:filtered-spectral-sequence}
设 \(\mathcal A\) 为交换范畴，\((C,F)\) 为 \(\mathcal A\) 中的递减滤过上链复形。取 \(E_0^{p,q}=F^pC^{p+q}/F^{p+1}C^{p+q}\)，并对 \(r\ge1\) 以\cref{b4:eq:spectral-page-formula}定义 \(E_r^{p,q}\)。原复形的微分诱导 \(\mathrm d_r:E_r^{p,q}\to E_r^{p+r,q-r+1}\)，其代表元公式为 \(\mathrm{d}_r[x]=[\mathrm{d}x]\)。这些页构成谱序列，且对滤过复形映射自然。
\end{theorem}
\begin{proof}
以下记 \(n=p+q\)。在模范畴中可以直接作代表元计算；一般情形按\cref{b4:lem:generalized-element-calculus}使用广义元素。记号 \(x\in Z\) 表示一个态射 \(x:T\to Z\)，\([x]\) 表示其经商态射所得的像。从商对象选代表、从像中选原像，或把子对象和中的态射写成两个分量之和时，都先沿相应满态射的拉回预合成。有限次操作可取共同的满态射细化，下面仍沿用原符号；同一等式内的态射具有相同的定义域，态射等式最终可消去预合成的满态射。

第零页微分的核在 \(F^pC^n\) 中的原像是 \(Z_1^{p,q}\)，边界的原像是 \(F^{p+1}C^n+\mathrm d(F^pC^{n-1})\)，所以 \(H(E_0,\mathrm d_0)=E_1\)。对 \(r\ge1\)，记第 \(r\) 页的分母为
\[
 L_r^{p,q}=Z_{r-1}^{p+1,q-1}
 +\mathrm d Z_{r-1}^{p-r+1,q+r-2}.
\]
由 \(\mathrm d^2=0\)，微分限制为 \(Z_r^{p,q}\to Z_r^{p+r,q-r+1}\)。它把分母第一项 \(u\) 送到目标分母的第二项 \(\mathrm d u\)，把第二项送到零。这些都是子对象之间的态射分解，故微分零化 \(L_r^{p,q}\) 到目标商对象的复合，唯一诱导真正的态射 \(\mathrm d_r\)，且 \(\mathrm d_r^2=0\)。

现比较 \(H(E_r,\mathrm{d}_r)\) 与 \(E_{r+1}\)。包含 \(Z_{r+1}^{p,q}\to Z_r^{p,q}\) 在第 \(r\) 页的像被 \(\mathrm d_r\) 零化，因而诱导态射 \(Z_{r+1}^{p,q}\to H(E_r)^{p,q}\)。为证它满，给定 \(\ker\mathrm d_r\) 的广义元素，先沿满态射 \(Z_r^{p,q}\to E_r^{p,q}\) 选代表 \(x\)。余圈条件说明 \(\mathrm d x\) 通过目标分母 \(L_r^{p+r,q-r+1}\) 分解。再沿到该分母的满态射
\[
 Z_{r-1}^{p+r+1,q-r}\oplus Z_{r-1}^{p+1,q-1}
 \longrightarrow L_r^{p+r,q-r+1},\qquad
 (u,v)\longmapsto u+\mathrm d v
\]
作拉回提升，便可在共同的满态射后写成
\[
 \mathrm{d}x=u+\mathrm{d}v,\qquad
 u\in Z_{r-1}^{p+r+1,q-r},\quad
 v\in Z_{r-1}^{p+1,q-1}.
\]
用 \(x-v\) 代替 \(x\) 不改变第 \(r\) 页的类，而且 \(\mathrm{d}(x-v)=u\) 通过 \(F^{p+r+1}\) 分解。拉回原像的泛性质遂使 \(x-v\) 通过 \(Z_{r+1}^{p,q}\) 分解。每个余圈都能在满态射后如此修正，再沿 \(\ker\mathrm d_r\to H(E_r)\) 的满态射选代表，便由广义元素判据知 \(Z_{r+1}^{p,q}\to H(E_r)^{p,q}\) 满。

再确定这个态射的核。设 \(x:T\to Z_{r+1}^{p,q}\) 在 \(H(E_r)\) 中的像为零，则它在 \(E_r\) 中的像落在 \(\mathrm d_r\) 的像内。先沿微分到其像的满态射提升，再沿来源页的商态射提升，得到 \(w\in Z_r^{p-r,q+r-1}\)，使 \([x]=[\mathrm d w]\)。因而 \(x-\mathrm d w\) 通过旧分母 \(L_r^{p,q}\) 分解；继续沿定义该和的直和到像的满态射提升，便得到
\[
 x=u+\mathrm{d}v+\mathrm{d}w,
\]
其中 \(u\in Z_{r-1}^{p+1,q-1}\)、\(v\in Z_{r-1}^{p-r+1,q+r-2}\)、\(w\in Z_r^{p-r,q+r-1}\)。若 \(x\in Z_{r+1}^{p,q}\)，则 \(\mathrm{d}u=\mathrm{d}x\in F^{p+r+1}\)，故 \(u\in Z_r^{p+1,q-1}\)。又 \(v+w\in F^{p-r}C^{n-1}\)，且 \(\mathrm{d}(v+w)=x-u\in F^pC^n\)，所以 \(v+w\in Z_r^{p-r,q+r-1}\)。因此核恰为
\[
 Z_r^{p+1,q-1}+\mathrm{d}Z_r^{p-r,q+r-1},
\]
正是第 \(r+1\) 页的分母。上述有限提升说明核中的任意广义元素在满态射后都来自这个和；反方向中，第一项在旧页已经为零，第二项给出旧页的边界。因此由广义元素正合性判据，核与所列子对象相等，取商即得 \(E_{r+1}^{p,q}\simeq H(E_r)^{p,q}\)。

滤过映射把以上子对象送入对应子对象，并与微分、商映射相容，所以诱导自然的页间映射。各页态射及同构由这些泛性质唯一确定，与计算时的满态射细化选择无关。
\end{proof}

构造的意义是逐步改进代表元：第 \(r\) 页的余圈不是已经满足 \(\mathrm{d}x=0\)，而是能够把 \(\mathrm{d}x\) 推到第 \(p+r\) 层。若 \(\mathrm{d}_r[x]\ne0\)，这种改进在此处遇到障碍；若它为零，则上一证明给出了下一次修正。

\ProseParagraphBreak
回到章首的四项双复形，\([a]\) 在第一页存留，而 \(\mathrm{d}a=c\) 仍落在中间列。用 \(-b\) 修正代表元后，\(\mathrm{d}(a-b)=-e\) 落入第二层，因而 \(\mathrm{d}_2[a]=-[e]\)；\secref{b4:sec:1405}会核对全部页与总复形。滤过谱序列的教材构造可对照 Weibel\cite[Section 5.4, Construction Theorem 5.4.1, Construction 5.4.6 and Lemma 5.4.7]{Weibel1994HomologicalAlgebra}；该书采用同调型指标，与本章固定的上同调型指标须作换算。
